Time-series forecasting is the task of learning the temporal dependencies that link past observations to future values in order to predict them. It is one of the tasks of time-series analysis, beside imputation, classification and anomaly detection, and it serves domains from energy and traffic to finance and weather, where a forecast feeds a decision rather than a report. Real series combine trend, seasonality and irregular fluctuation, and are often non-stationary, so a forecaster has to separate the structure worth extrapolating from the noise around it\autocite{wangDeepTimeSeries2026}.

\subsection*{Classic Aliasing}
Predictive fidelity is bounded by the digital representation of the signal, whose governing limit is the Nyquist-Shannon sampling theorem\autocite{shannonCommunicationPresenceNoise1949,oppenheimDiscreteTimeSignalProcessing2010}: a band-limited continuous-time signal is perfectly reconstructible only if sampled uniformly at a rate $f_s$ strictly greater than twice its highest frequency component $f_{\max}$, that is $f_s>2f_{\max}$.

The threshold $f_N = f_s/2$ is the Nyquist frequency. When the criterion is violated, components above $f_N$ fold back into the lower spectrum and become indistinguishable from genuine low-frequency content, an irreversible loss known as aliasing, whereas satisfying it is taken as sufficient for a faithful discrete representation.

\subsection*{Evolution of Forecasting Methodologies}
Classical statistical models describe a series through an explicit equation with few parameters, estimated from that series alone. The next value is a weighted combination of past values, past errors or smoothed components such as level, trend and season, as in Exponential Smoothing and autoregressive integrated moving average (ARIMA) models\autocite[chs.~8--9]{hyndmanForecastingPrinciplesPractice2021}. They are interpretable and need little data, but assume linear dynamics and a stationary series. Deep models relax these assumptions by learning the dependencies from data, first with recurrent and convolutional networks and then with Transformers, which capture long-range dependence at a cost quadratic in sequence length\autocite{vaswaniAttentionAllYou2017,wangDeepTimeSeries2026}. Patch-based architectures answer that cost by grouping neighbouring observations into tokens\autocite{nieTimeSeriesWorth2023}. Chronos quantises values into discrete tokens\autocite{ansariChronosLearningLanguage2024}; Chronos-Bolt replaces quantisation with patching and a direct multi-quantile head\autocite{ansariChronosLearningLanguage2024a,FastAccurateZeroshot2024a,AmazonChronosbolttinyHugging2025}.

Patch-based tokenisation has become the default input stage of time-series foundation models. Its assumed price is resolution, with the sampling theorem deciding what survives; this report argues that the price is of a different kind. Where the period of a signal is commensurate with the patch geometry, patches repeat, and a frequency the sampler represented perfectly may become poorly observable to the model.

The distinction matters wherever such a model sits in a monitoring or control loop. A blind spot that followed from the configuration rather than from the data would be invisible to an operator who never derived it and predictable to anyone who did. It would be an engineering defect and a security property at once, and is treated here as a candidate for both.

\subsection*{Application Domain and Agent Objectives}
The reference setting is a designed scenario: a photovoltaic array supplying an internal load, a local storage battery and the external grid, with a control agent that balances load, dispatches the battery and detects anomalies on short-horizon forecasts. Measured telemetry comes from the SolarTech Lab dataset\autocite{levaPhotovoltaicPowerWeather2020}; the concepts, states and relations of the scenario are formalised as an ontology in the Web Ontology Language (OWL2), and the agent bases its monitoring and power-routing decisions on forecasts from the Chronos-Bolt family.

\subsection*{Scope}
The target architecture is Chronos-Bolt-Tiny, with default patch size $P=16$, stride $S=16$, context window $T=2048$ and a direct quantile head\autocite{AmazonChronosbolttinyHugging2025,FastAccurateZeroshot2024a}. The objective is to test the assumption that such a model preserves the frequency content of Nyquist-compliant signals, to characterise how $f_s$, $P$ and $S$ interact in determining when it does not. It is also explored how the patch geometry affects the representation of complex inputs, from synthetic signals to photovoltaic telemetry. This study excludes geometries with $S>P$, other foundation models and live attacks on a deployed system.

\subsection*{Report Structure}
\autoref{sec:tokenisation} models patch tokenisation as an operator and derives the phase-locking conditions and the candidate set they define. \autoref{sec:agent} presents the ontology and the control agent, states the hypotheses to be tested and the proposed mitigation. \autoref{sec:signals} describes the synthetic and measured signals, and \autoref{sec:methodology} the four experimental steps, the Bayesian models and the pre-registered decision rules. \autoref{sec:results} reports the frequency sweep, the decomposition of generated signals, the forecasts on the solar telemetry and the Bayesian analysis. \autoref{sec:discussion} reads the Bayesian claims one by one, evaluates the mitigation and collects the limitations of the study, \autoref{sec:security} assesses the risk under the Artificial Intelligence Risk Management Framework, and \autoref{sec:conclusions} ties the conclusions. Eight appendices give the detail behind each step: the ontology (\autoref{sec:appOntology}), the retraining recipe (\autoref{sec:appRetraining}), the two signal generators (\autoref{sec:appTSMixup} and \autoref{sec:appKernelSynth}), the Bayesian specifications (\autoref{sec:appBayes}), and the sweep, decomposition and photovoltaic analyses (\autoref{sec:appSweep}, \autoref{sec:appDecomp} and \autoref{sec:appSolar}).
